\begin{thebibliography}{99}

\bibitem{BaierBansal}
S. Baier and A. Bansal,
\emph{The large sieve with power moduli for $\mathbb Z[i]$},
International Journal of Number Theory \textbf{14} (2018), no. 10,
2737--2756.
\href{https://doi.org/10.1142/S1793042118501658}{doi:10.1142/S1793042118501658};
\href{https://arxiv.org/abs/1802.08964}{arXiv:1802.08964}.

\bibitem{BIMS}
V. Bhargava, G. Ivanyos, R. Mittal, and N. Saxena,
\emph{Irreducibility and deterministic $r$-th root finding over finite fields},
Proceedings of the 2017 ACM International Symposium on Symbolic and
Algebraic Computation (ISSAC 2017), ACM, 2017, 37--44.
\href{https://doi.org/10.1145/3087604.3087620}{doi:10.1145/3087604.3087620};
\href{https://www.cse.iitk.ac.in/users/nitin/papers/irred-and-rth-root.pdf}{author manuscript}.

\bibitem{BCSS}
K. Billington, M. Cheng, J. Schettler, and A. I. Suriajaya,
\emph{The average number of Goldbach representations and zero-free regions
of the Riemann zeta-function},
International Journal of Number Theory \textbf{21} (2025), no. 2, 289--316.
\href{https://doi.org/10.1142/S1793042125500150}{doi:10.1142/S1793042125500150};
\href{https://arxiv.org/abs/2306.09102}{arXiv:2306.09102}.

\bibitem{BGL}
V. Blomer, L. Goldmakher, and B. Louvel,
\emph{$L$-functions with $n$-th-order twists},
International Mathematics Research Notices \textbf{2014} (2014), no. 7,
1925--1955.
\href{https://doi.org/10.1093/imrn/rns257}{doi:10.1093/imrn/rns257};
\href{https://arxiv.org/abs/1112.1650}{arXiv:1112.1650}.

\bibitem{Clay}
Clay Mathematics Institute,
\emph{Riemann Hypothesis}, Millennium Prize Problems,
\url{https://www.claymath.org/millennium/riemann-hypothesis/}.

\bibitem{DFDS}
C. David, A. de Faveri, A. Dunn, and J. Stucky,
\emph{Non-vanishing for cubic Hecke $L$-functions},
preprint, arXiv:2410.03048v2 (2026).
\href{https://arxiv.org/abs/2410.03048}{arXiv:2410.03048}.

\bibitem{Dirichlet}
P. G. L. Dirichlet,
\emph{Beweis des Satzes, da\ss{} jede unbegrenzte arithmetische
Progression, deren erstes Glied und Differenz ganze Zahlen ohne
gemeinschaftlichen Factor sind, unendlich viele Primzahlen enth\"alt},
Abhandlungen der K\"oniglichen Preussischen Akademie der Wissenschaften
zu Berlin (1837), 45--81.
English translation by R. Stephan,
\href{https://arxiv.org/abs/0808.1408}{arXiv:0808.1408}.

\bibitem{DR}
A. Dunn and M. Radziwi\l\l,
\emph{Bias in cubic Gauss sums: Patterson's conjecture},
Annals of Mathematics \textbf{200} (2024), no. 3, 967--1057.
\href{https://doi.org/10.4007/annals.2024.200.3.3}{doi:10.4007/annals.2024.200.3.3};
\href{https://arxiv.org/abs/2109.07463}{arXiv:2109.07463}.

\bibitem{FG}
J. B. Friedlander and D. A. Goldston,
\emph{Sums of three or more primes},
Transactions of the American Mathematical Society \textbf{349} (1997),
no. 1, 287--310.
\href{https://doi.org/10.1090/S0002-9947-97-01652-8}{doi:10.1090/S0002-9947-97-01652-8}.

\bibitem{Galbraith}
S. D. Galbraith,
\emph{Mathematics of Public Key Cryptography},
Cambridge University Press, 2012.
\href{https://doi.org/10.1017/CBO9781139012843}{doi:10.1017/CBO9781139012843}.
Section and algorithm references use the author's extended and corrected
version 2.0, \href{https://www.math.auckland.ac.nz/~sgal018/crypto-book/ch2.pdf}{Chapter 2}.

\bibitem{GZ}
P. Gao and L. Zhao,
\emph{Moments of quadratic Hecke $L$-functions of imaginary quadratic
number fields},
Journal of Number Theory \textbf{209} (2020), 359--377.
\href{https://doi.org/10.1016/j.jnt.2019.09.002}{doi:10.1016/j.jnt.2019.09.002};
\href{https://arxiv.org/abs/1707.00091}{arXiv:1707.00091}.

\bibitem{GZSextic}
P. Gao and L. Zhao,
\emph{Moments and one level density of sextic Hecke $L$-functions},
Functiones et Approximatio Commentarii Mathematici \textbf{70} (2024),
no. 1, 7--28.
\href{https://doi.org/10.7169/facm/2057}{doi:10.7169/facm/2057};
\href{https://arxiv.org/abs/2201.01885}{arXiv:2201.01885}.

\bibitem{GL}
L. Goldmakher and B. Louvel,
\emph{A quadratic large sieve inequality over number fields},
Mathematical Proceedings of the Cambridge Philosophical Society
\textbf{154} (2013), no. 2, 193--212.
\href{https://doi.org/10.1017/S0305004112000370}{doi:10.1017/S0305004112000370};
\href{https://arxiv.org/abs/1112.1642}{arXiv:1112.1642}.

\bibitem{GM}
L. Guth and J. Maynard,
\emph{New large value estimates for Dirichlet polynomials},
Annals of Mathematics \textbf{203} (2026), no. 2, 623--675.
\href{https://doi.org/10.4007/annals.2026.203.2.6}{doi:10.4007/annals.2026.203.2.6};
\href{https://arxiv.org/abs/2405.20552}{arXiv:2405.20552}.

\bibitem{Hadamard}
J. Hadamard,
\emph{Sur la distribution des z\'eros de la fonction $\zeta(s)$ et ses
cons\'equences arithm\'etiques},
Bulletin de la Soci\'et\'e Math\'ematique de France \textbf{24} (1896),
199--220.
\href{https://doi.org/10.24033/bsmf.545}{doi:10.24033/bsmf.545}.

\bibitem{HBQuadratic}
D. R. Heath-Brown,
\emph{A mean value estimate for real character sums},
Acta Arithmetica \textbf{72} (1995), no. 3, 235--275.
\href{https://doi.org/10.4064/aa-72-3-235-275}{doi:10.4064/aa-72-3-235-275}.

\bibitem{HB}
D. R. Heath-Brown,
\emph{Kummer's conjecture for cubic Gauss sums},
Israel Journal of Mathematics \textbf{120}, part A (2000), 97--124.
\href{https://doi.org/10.1007/s11856-000-1273-y}{doi:10.1007/s11856-000-1273-y}.

\bibitem{Hecke}
E. Hecke,
\emph{Eine neue Art von Zetafunktionen und ihre Beziehungen zur
Verteilung der Primzahlen. Erste Mitteilung},
Mathematische Zeitschrift \textbf{1} (1918), 357--376;
\emph{Zweite Mitteilung}, Mathematische Zeitschrift \textbf{6} (1920),
11--51.
\href{https://doi.org/10.1007/BF01465095}{doi:10.1007/BF01465095}
(first paper).

\bibitem{Huxley}
M. N. Huxley,
\emph{The large sieve inequality for algebraic number fields},
Mathematika \textbf{15} (1968), no. 2, 178--187.
\href{https://doi.org/10.1112/S0025579300002540}{doi:10.1112/S0025579300002540}.

\bibitem{Kubota}
T. Kubota,
\emph{On automorphic functions and the reciprocity law in a number field},
Lectures in Mathematics, Department of Mathematics, Kyoto University,
no. 2, Kinokuniya Book-Store, Tokyo, 1969.
\url{http://hdl.handle.net/2433/84907}.

\bibitem{MP}
J. Maynard and K. Pratt,
\emph{Half-isolated zeros and zero-density estimates},
International Mathematics Research Notices \textbf{2024} (2024), no. 19,
12978--13014.
\href{https://doi.org/10.1093/imrn/rnae191}{doi:10.1093/imrn/rnae191};
\href{https://arxiv.org/abs/2206.11729}{arXiv:2206.11729}.

\bibitem{MilneCFT}
J. S. Milne,
\emph{Class Field Theory}, course notes, version 4.03, 6 August 2020.
\url{https://www.jmilne.org/math/CourseNotes/CFT.pdf}.

\bibitem{OpenAICubicFirstMoment}
OpenAI,
\emph{An unconditional first moment for cubic Gauss sums},
OpenAI Math Release preprint
\href{https://github.com/openai/math/blob/main/preprints/An-unconditional-first-moment-for-cubic-Gauss-sums-September-25-2026/paper.pdf}{OAI:An-unconditional-first-moment-for-cubic-Gauss-sums-September-25-2026},
2026.

\bibitem{Patterson}
S. J. Patterson,
\emph{A cubic analogue of the theta series},
Journal f\"ur die reine und angewandte Mathematik \textbf{296} (1977),
125--161.
\href{https://doi.org/10.1515/crll.1977.296.125}{doi:10.1515/crll.1977.296.125}.

\bibitem{Riemann}
B. Riemann,
\emph{Ueber die Anzahl der Primzahlen unter einer gegebenen Gr\"osse},
Monatsberichte der K\"oniglichen Preussischen Akademie der Wissenschaften
zu Berlin (1859), 671--680.
Manuscript, transcription, and translation at
\url{https://www.claymath.org/library/historical/riemann/r.html}.

\bibitem{TaoNonresidues}
T. Tao,
\emph{The Elliott--Halberstam conjecture implies the Vinogradov least
quadratic nonresidue conjecture},
Algebra \& Number Theory \textbf{9} (2015), no. 4, 1005--1034.
\href{https://doi.org/10.2140/ant.2015.9.1005}{doi:10.2140/ant.2015.9.1005};
\href{https://msp.org/ant/2015/9-4/ant-v9-n4-p09-p.pdf}{publisher PDF}.

\bibitem{TZ}
J. Thorner and A. Zaman,
\emph{A unified and improved Chebotarev density theorem},
Algebra \& Number Theory \textbf{13} (2019), no. 5, 1039--1068.
\href{https://doi.org/10.2140/ant.2019.13.1039}{doi:10.2140/ant.2019.13.1039};
\href{https://arxiv.org/abs/1803.02823}{arXiv:1803.02823}.

\bibitem{ValleePoussin}
C.-J. de la Vall\'ee Poussin,
\emph{Recherches analytiques sur la th\'eorie des nombres premiers.
Premi\`ere partie: La fonction $\zeta(s)$ de Riemann et les nombres
premiers en g\'en\'eral},
Annales de la Soci\'et\'e scientifique de Bruxelles \textbf{20},
deuxi\`eme partie (1896), 183--256.
Reprinted in \emph{Recherches analytiques sur la th\'eorie des nombres
premiers}, Hayez, Brussels, 1897,
\url{https://books.google.com/books?id=7e0GAAAAYAAJ}.

\end{thebibliography}
\end{document}

